Much of modern statistics and machine learning can be phrased as choosing a parameter vector $x\in\R^d$ to minimize an objective:
\begin{equation}
  \min_{x \in \R^d} f(x),
  \label{eq:unconstrained-min}
\end{equation}
This unconstrained problem is the core primitive behind many later constructions: penalties, constraints, stochastic approximations, and even some sampling algorithms are all variations on it.
Many statistical procedures reduce to minimizing a loss plus regularization; see \citet[\S9.1]{boyd_vandenberghe_2004_convex}.

Local analysis begins with a concrete question: if we perturb the current point $x$ in direction $v$, what is the first-order change in $f$?
The differential $Df(x)[v]$ answers exactly that question, and it is linear in the direction $v$.
That linearity is what makes the gradient useful: one vector can summarize all directional slopes at once.
In Euclidean coordinates every linear functional on $\R^d$ can be written uniquely in the form
\[
  v \mapsto \ip{a,v}
  \qquad\text{for some } a\in\R^d \text{ and all } v\in\R^d,
\]
so we represent first-order information by a vector.
For a differentiable function $f$, the gradient is defined by the identity
\[
  Df(x)[v]=\ip{\nabla f(x),v}
  \qquad\text{for all } v\in\R^d.
\]
Thus the gradient is not an extra construction; it is simply the Euclidean way of writing the linear approximation of the objective.
For example, if $f(x)=a^\top x$ then $\nabla f(x)=a$, while for $f(x)=\frac12\norm{x}_2^2$ we have $\nabla f(x)=x$.
For now we stay in the unconstrained setting and develop the derivative calculus itself.
The multiplier language belongs to constrained optimization, where the new issue is feasibility rather than local change in $f$.
In convex optimization this first-order information becomes especially powerful because the local model controls the function globally.

We will use regularized logistic regression as a running example throughout the optimization chapter.
It is simple enough to differentiate explicitly, but rich enough to show how curvature, conditioning, and regularization interact.

\begin{example}[Logistic regression (with ridge regularization)]
\label{ex:logistic-reg}
Let data $(x_i,y_i)_{i=1}^n$ satisfy $x_i \in \R^d$ and $y_i \in \{0,1\}$.
In logistic regression we model
\begin{equation}
  y_i \mid x_i \sim \mathrm{Bernoulli}(p_i(\theta)),
  \qquad
  p_i(\theta) = \sigma(x_i^\top \theta)
  = \frac{1}{1+e^{-x_i^\top \theta}}.
  \label{eq:logistic-model}
\end{equation}
The log-likelihood is
\begin{align}
  \ell(\theta)
  &= \log p(y_1,\dots,y_n \mid x_1,\dots,x_n,\theta) \nonumber \\
  &= \sum_{i=1}^n \Bigl[y_i x_i^\top \theta - \log\bigl(1+e^{x_i^\top \theta}\bigr)\Bigr].
  \label{eq:logistic-loglik}
\end{align}
Equivalently, the (convex) negative log-likelihood is
\begin{equation}
  L(\theta) = -\ell(\theta)
  = \sum_{i=1}^n \Bigl[\log\bigl(1+e^{x_i^\top \theta}\bigr) - y_i x_i^\top \theta\Bigr].
  \label{eq:logistic-nll}
\end{equation}
A common regularized objective (ridge) is
\begin{equation}
  s(\theta) = L(\theta) + \frac{m}{2}\norm{\theta}_2^2,
  \qquad m>0.
  \label{eq:logistic-ridge-objective}
\end{equation}
\end{example}

The objective \eqref{eq:logistic-ridge-objective} is a clean example of a smooth convex optimization problem with a finite-sum structure (a sum over data points).
It also lets us see, explicitly, how curvature and conditioning depend on the data matrix $X$ and on regularization.

\begin{lemma}[Gradient and Hessian for logistic regression]
\label{lem:logistic-grad-hess}
Let $X \in \R^{n \times d}$ have rows $x_i^\top$, and let $p(\theta) \in \R^n$ have entries $p_i(\theta)=\sigma(x_i^\top\theta)$.
Then
\begin{equation}
  \nabla L(\theta) = X^\top \bigl(p(\theta) - y\bigr),
  \qquad
  \Hess L(\theta) = X^\top W(\theta) X,
  \label{eq:logistic-grad-hess}
\end{equation}
where $W(\theta) = \mathrm{diag}(p_i(\theta)(1-p_i(\theta)))$.
Moreover, since $0 < p(1-p) \le 1/4$ for all $p \in (0,1)$,
\begin{equation}
  0 \preceq W(\theta) \preceq \tfrac14 I
  \quad\Rightarrow\quad
  0 \preceq X^\top W(\theta) X \preceq \tfrac14 X^\top X.
  \label{eq:logistic-hess-bound}
\end{equation}
For the ridge objective $s(\theta)=L(\theta)+\tfrac{m}{2}\norm{\theta}_2^2$,
\begin{equation}
  \Hess s(\theta) = X^\top W(\theta) X + m I.
  \label{eq:logistic-ridge-hess}
\end{equation}
\end{lemma}
\begin{proof}
Write $z_i(\theta)=x_i^\top\theta$.
The derivative of $\log(1+e^{z})$ is $\frac{e^z}{1+e^z}=\sigma(z)$, so differentiating \eqref{eq:logistic-nll} gives
\[
\nabla L(\theta)
= \sum_{i=1}^n \bigl(\sigma(z_i(\theta)) - y_i\bigr)x_i
= X^\top(p(\theta)-y).
\]
Since $\sigma'(z)=\sigma(z)(1-\sigma(z))$, another derivative yields
\[
\Hess L(\theta)=\sum_{i=1}^n p_i(\theta)(1-p_i(\theta))\,x_i x_i^\top = X^\top W(\theta)X.
\]
The bound $p(1-p)\le 1/4$ is maximized at $p=1/2$, which gives \eqref{eq:logistic-hess-bound}.
Finally, $\Hess\left(\tfrac{m}{2}\norm{\theta}_2^2\right)=mI$, so \eqref{eq:logistic-ridge-hess} follows by additivity.
\end{proof}

The Hessian for logistic regression is a weighted Gram matrix $X^\top W X$.
When a datum is ``confident'' (i.e., $p_i(\theta)$ close to $0$ or $1$), its weight $p_i(1-p_i)$ is small, so it contributes little curvature; the maximum curvature occurs near the decision boundary ($p_i \approx 1/2$).
Adding ridge regularization shifts all eigenvalues up by $m$, which improves conditioning and helps algorithms like gradient descent and Newton's method behave more predictably.
